% ===================================================================
%  AntennaForge — Antenna Pattern Fidelity Assessment
%  Per-Antenna Grading, Shortcomings, and Potential Improvements
% ===================================================================
\documentclass[11pt,a4paper]{article}

% ── Packages ──────────────────────────────────────────────────────
\usepackage[utf8]{inputenc}
\usepackage[T1]{fontenc}
\usepackage{lmodern}
\usepackage[margin=1in]{geometry}
\usepackage{xcolor}
\usepackage{hyperref}
\usepackage{booktabs}
\usepackage{longtable}
\usepackage{amsmath,amssymb}
\usepackage{titlesec}
\usepackage{fancyhdr}
\usepackage{parskip}
\usepackage{enumitem}
\DeclareMathOperator{\sinc}{sinc}

% ── Colour Definitions ───────────────────────────────────────────
\definecolor{afblue}{HTML}{3B82F6}
\definecolor{afdark}{HTML}{0F172A}
\definecolor{afgray}{HTML}{475569}
\definecolor{afred}{HTML}{DC2626}
\definecolor{afgreen}{HTML}{16A34A}
\definecolor{aforange}{HTML}{EA580C}

% ── Hyperref Setup ───────────────────────────────────────────────
\hypersetup{
    colorlinks=true,
    linkcolor=afblue,
    urlcolor=afblue,
    pdftitle={AntennaForge — Antenna Pattern Fidelity Assessment},
    pdfauthor={AntennaForge Project},
}

% ── Header / Footer ─────────────────────────────────────────────
\pagestyle{fancy}
\fancyhf{}
\fancyhead[L]{\bfseries AntennaForge — Pattern Fidelity Assessment}
\setlength{\headheight}{13.6pt}
\fancyhead[R]{\thepage}
\renewcommand{\headrulewidth}{0.4pt}

% ── Section Formatting ───────────────────────────────────────────
\titleformat{\section}{\Large\bfseries\color{afdark}}{}{0em}{\thesection\quad}
\titleformat{\subsection}{\large\bfseries\color{afblue}}{}{0em}{\thesubsection\quad}
\titleformat{\subsubsection}{\normalsize\bfseries\color{afgray}}{}{0em}{\thesubsubsection\quad}

% ── Grade box command ────────────────────────────────────────────
\newcommand{\grade}[2]{%
    \par\noindent\fcolorbox{#1}{#1!10}{%
        \parbox{0.95\textwidth}{\textbf{\color{#1}Grade: #2}}%
    }\par\medskip
}

\begin{document}

% ── Title ────────────────────────────────────────────────────────
\begin{center}
    {\Huge\bfseries\color{afdark} Antenna Pattern Fidelity Assessment}\\[0.5em]
    {\large\color{afgray} Per-Antenna Grading as Parameterised Synthetic Models}\\[1em]
    {\color{afgray} AntennaForge Project \quad---\quad March 2026}
\end{center}
\vspace{1em}

% ── Preamble ─────────────────────────────────────────────────────
\section{Scope and Grading Criteria}

This document assesses each AntennaForge antenna type on its merits as a
\textbf{parameterised synthetic pattern generator}---not against a full-wave
electromagnetic solver such as NEC, FEKO, or HFSS. The grading criteria are:

\begin{itemize}[nosep]
    \item \textbf{Physical basis} --- Is the front-lobe model grounded in
          aperture theory, diffraction, or analytical EM?
    \item \textbf{Parameter utilisation} --- Are all inputs to
          \texttt{compute\_point\_gain()} meaningfully used?
    \item \textbf{Frequency behaviour} --- Does the pattern evolve correctly
          with frequency without manual intervention?
    \item \textbf{Back hemisphere} --- Is the back-lobe model reasonable and
          artefact-free?
    \item \textbf{Extensibility} --- How easily can the model be improved
          without restructuring?
\end{itemize}

\noindent
Grades use a standard letter scale: \textbf{A+} (excellent),
\textbf{A} (strong), \textbf{A$-$}, \textbf{B+}, \textbf{B},
\textbf{B$-$}, \textbf{C+}, etc.

\bigskip
\noindent
\textit{Source files assessed:}
\texttt{antennas/monopole/monopole.py},
\texttt{antennas/dish/dish.py},
\texttt{antennas/panel/panel.py},
\texttt{antennas/horn/horn.py},
\texttt{antennas/array/array\_antenna.py},
\texttt{antennas/lpda/lpda.py},
\texttt{core/pattern\_math.py}.


% =================================================================
\newpage
\section{Monopole / Dipole}
% =================================================================

\grade{afgreen}{A+}

\subsection{Model Summary}

The monopole uses the \textbf{exact analytical dipole formula}:
\[
G(\theta) = G_\text{peak} \cdot
    \left[\frac{\cos\!\bigl(\tfrac{kL}{2}\sin\theta\bigr) - \cos\!\bigl(\tfrac{kL}{2}\bigr)}
         {1 - \cos\!\bigl(\tfrac{kL}{2}\bigr)}\right]^{2}
    \cdot \frac{1}{\cos^2\theta}
\]
where $kL = 2\pi L$ and $L$ is the element length in wavelengths.

\subsection{Parameter Assessment}

\begin{longtable}{@{}llp{6cm}@{}}
\toprule
\textbf{Parameter} & \textbf{Grade} & \textbf{Notes} \\
\midrule
\endhead
\texttt{element\_length\_wavelengths} & A+ & Drives the entire pattern shape. Quarter-wave, half-wave, 5/8-wave all produce correct, distinct patterns. \\
\texttt{monopole\_physical\_length\_m} & A+ & Enables automatic frequency scaling: $L(f) = \text{phys} / \lambda(f)$. Pattern evolves across band exactly as a fixed physical element would. \\
\texttt{gain\_peak} & A & Passed through. The only user knob---everything else is physics-derived. \\
Azimuth pattern & A+ & Omnidirectional (\texttt{has\_az\_pattern = False}). Correct for a vertical element over a ground plane. \\
Degenerate handling & A & Full-wave ($L = 1.0\lambda$) singularity caught; falls back to $\cos^n$ model. \\
\bottomrule
\end{longtable}

\subsection{Shortcomings}

\begin{enumerate}[nosep]
    \item \textbf{No ground-plane size effects.} A real monopole's pattern depends on the finite ground plane radius; the model assumes an infinite ground plane.
    \item \textbf{Gain not derived from directivity.} \texttt{gain\_peak} is a user knob rather than computed from pattern integration ($D = 1.64$ for a quarter-wave monopole, $2.15$ for a half-wave dipole).
    \item \textbf{No impedance information.} Feed-point impedance ($36.5 + j21.25\,\Omega$ for a quarter-wave monopole) is not computed.
    \item \textbf{Full-wave fallback is approximate.} The $\cos^n$ fallback at $L = 1.0\lambda$ does not match the true pattern (a four-lobed cloverleaf).
\end{enumerate}

\subsection{Potential Improvements}

\begin{enumerate}[nosep]
    \item Compute theoretical directivity from element length and auto-set \texttt{gain\_peak} when the user does not override it.
    \item Add finite ground-plane correction using image theory with a truncation factor.
    \item Compute feed-point impedance from Hallen's or King's equations for thin-wire dipoles.
    \item Replace the full-wave $\cos^n$ fallback with the proper limit of the dipole formula (L'H\^opital's rule at $\cos(kL/2) \to 1$).
\end{enumerate}


% =================================================================
\newpage
\section{Dish (Parabolic Reflector)}
% =================================================================

\grade{afgreen}{A}

\subsection{Model Summary}

Circular aperture model using the \textbf{Jinc-squared (Airy) pattern}:
\[
G(\theta) = G_\text{peak} \cdot \eta \cdot
    \left[\frac{2\,J_1(u)}{u}\right]^{2} \cdot T(r)
\]
where $u = 1.6163 \cdot \sin\theta / \sin(\theta_{3\text{dB}}/2)$,
$\eta$ is aperture efficiency, and $T(r)$ is a Gaussian feed taper envelope.

The constant $1.6163$ is the exact solution to $[2J_1(K)/K]^2 = 0.5$,
guaranteeing $-3$\,dB at the stated beamwidth.

\subsection{Parameter Assessment}

\begin{longtable}{@{}llp{6cm}@{}}
\toprule
\textbf{Parameter} & \textbf{Grade} & \textbf{Notes} \\
\midrule
\endhead
Front lobe (Jinc$^2$) & A+ & Exact circular aperture diffraction. Sidelobe positions and amplitudes match theory ($-17.6$\,dB first sidelobe for uniform illumination). \\
\texttt{dish\_efficiency} & B+ & Realistic knob (0.55--0.65 typical), but not decomposed into spillover, blockage, and surface-error components. \\
\texttt{feed\_edge\_taper\_db} & A & Gaussian envelope $\exp(-r^2/2\sigma^2)$ correctly suppresses sidelobes; $\sigma$ derived from taper dB at dish edge. \\
\texttt{ftb\_ratio\_db} & B & User knob. A real dish FTB depends on edge diffraction, feed spillover, and strut scattering. \\
Bessel $J_1$ & A & Uses \texttt{scipy.special.j1} when available; Abramowitz \& Stegun polynomial fallback accurate to $\sim 8$ significant figures. \\
Back lobe & B & Heuristic $(1 - \cos\phi)/2$ weighting with $|\cos\theta|^{0.3n}$ elevation taper. Smooth, artefact-free, but not physics-derived. \\
\texttt{sigmoid\_k} & B & Default $k = 40$ gives $\sim\!5^\circ$ transition at $\pm 90^\circ$; sharper than physical reality for a dish. \\
\bottomrule
\end{longtable}

\subsection{Shortcomings}

\begin{enumerate}[nosep]
    \item \textbf{No frequency-dependent beamwidth.} A real dish has $\text{BW} \propto \lambda / D$. Beamwidth is static unless the engine's frequency-dependent BW feature is manually enabled.
    \item \textbf{Efficiency is a single scalar.} Real aperture efficiency decomposes into illumination efficiency, spillover, blockage ratio, and surface RMS error---each frequency-dependent.
    \item \textbf{No aperture geometry input.} Dish diameter ($D$), focal length ($f$), and $f/D$ ratio are not modelled. Beamwidth and gain should be derivable from $D$ and $\lambda$.
    \item \textbf{Back lobe model is generic.} Identical structure to all other antenna types; does not model feed spillover past the dish rim.
    \item \textbf{Strut diffraction not modelled.} Quadrupod/tripod support struts create characteristic sidelobe spikes at specific angles.
\end{enumerate}

\subsection{Potential Improvements}

\begin{enumerate}[nosep]
    \item Accept dish diameter as an input; derive $\text{BW} = 70\lambda / D$ and $G = \eta (\pi D / \lambda)^2$ automatically.
    \item Decompose efficiency into sub-components: illumination ($\eta_\text{ill}$), spillover ($\eta_\text{sp}$), blockage ($\eta_\text{bl}$), surface ($\eta_\text{surf} = \exp[-(4\pi\epsilon/\lambda)^2]$ Ruze equation).
    \item Model feed spillover as the back-lobe contribution rather than using the generic $(1-\cos\phi)/2$ formula.
    \item Add optional strut diffraction spikes at user-specified angles and widths.
\end{enumerate}


% =================================================================
\newpage
\section{Panel (Sector Antenna)}
% =================================================================

\grade{afgreen}{A}

\subsection{Model Summary}

Rectangular aperture model using \textbf{separable sinc-squared} pattern:
\[
G(\phi, \theta) = G_\text{peak} \cdot \sinc^2(u_\text{az}) \cdot \sinc^2(u_\text{el})
\]
where $u = 1.3916 \cdot \sin\theta / \sin(\theta_{3\text{dB}}/2)$ and
$\sinc(x) = \sin(x)/x$. The constant $1.3916$ solves $\sinc^2(K) = 0.5$.

\subsection{Parameter Assessment}

\begin{longtable}{@{}llp{6cm}@{}}
\toprule
\textbf{Parameter} & \textbf{Grade} & \textbf{Notes} \\
\midrule
\endhead
Front lobe ($\sinc^2$) & A+ & Exact rectangular aperture diffraction. First sidelobe at $-13.2$\,dB matches uniform illumination theory. Separable in az/el---valid for a rectangular aperture. \\
\texttt{mechanical\_tilt\_deg} & A & Clean elevation offset $\theta_\text{eff} = \theta - \text{tilt}$. Models real base-station mechanical downtilt. \\
\texttt{ftb\_ratio\_db} & B & User knob. Real panel FTB depends on reflector depth and edge treatment. \\
Back lobe & B & Same generic heuristic as all antenna types. \\
\texttt{sigmoid\_k} & B & Default $k = 40$ is too sharp for a panel antenna's gradual rolloff past $\pm 90^\circ$. \\
\texttt{n\_az}, \texttt{n\_el} & C+ & Only \texttt{n\_el} is used (for back-lobe broadening); \texttt{n\_az} is dead. \\
\bottomrule
\end{longtable}

\subsection{Shortcomings}

\begin{enumerate}[nosep]
    \item \textbf{No electrical downtilt.} Real sector antennas use phase-shifted column feeds to steer the beam in elevation without mechanical tilt. Only mechanical tilt is modelled.
    \item \textbf{Uniform illumination only.} The $\sinc^2$ pattern corresponds to uniform aperture illumination. Real panels use tapered illumination (e.g.\ Taylor, Dolph--Chebyshev) to control sidelobes, producing a wider main beam with lower sidelobes.
    \item \textbf{No aperture dimensions.} Panel width/height are not inputs; beamwidth is a user knob rather than derived from aperture size.
    \item \textbf{No frequency-dependent beamwidth.} BW should scale as $\lambda / W$ where $W$ is aperture width.
    \item \textbf{\texttt{n\_az} is unused.} Passed into \texttt{compute\_point\_gain()} but never referenced---wasted parameter.
\end{enumerate}

\subsection{Potential Improvements}

\begin{enumerate}[nosep]
    \item Add electrical downtilt via progressive phase across elevation elements.
    \item Accept aperture dimensions (width, height); derive beamwidths as $\text{BW} \approx 51\lambda / W$ (for $-13$\,dB first SLL) or $70\lambda / W$ (uniform).
    \item Support illumination taper options (uniform, Taylor, Chebyshev) that modify the pattern from pure $\sinc^2$ to the appropriate windowed shape.
    \item Remove or utilise the \texttt{n\_az} parameter.
\end{enumerate}


% =================================================================
\newpage
\section{Horn (Pyramidal)}
% =================================================================

\grade{afgreen}{A}

\subsection{Model Summary}

Standard pyramidal horn model with \textbf{asymmetric E/H-plane patterns}:
\begin{align}
    \text{E-plane (elevation):}\quad & G_\text{el} = \sinc^2(u_\text{el})
        & u_\text{el} &= 1.3916 \cdot \frac{\sin\theta_\text{el}}{\sin(\theta_{\text{el},3\text{dB}}/2)} \\
    \text{H-plane (azimuth):}\quad & G_\text{az} = \left[\frac{\cos v}{1 - (2v/\pi)^2}\right]^{2}
        & v &= 1.1891 \cdot \frac{\sin\theta_\text{az}}{\sin(\theta_{\text{az},3\text{dB}}/2)}
\end{align}
Total front-lobe gain: $G = G_\text{peak} \cdot G_\text{az} \cdot G_\text{el}$ (Pattern Multiplication Theorem).

\subsection{Parameter Assessment}

\begin{longtable}{@{}llp{6cm}@{}}
\toprule
\textbf{Parameter} & \textbf{Grade} & \textbf{Notes} \\
\midrule
\endhead
E-plane ($\sinc^2$) & A & Uniform illumination model. Correct for the electric-field distribution in the E-plane of a pyramidal horn. \\
H-plane (cosine taper) & A+ & The $\cos v / [1-(2v/\pi)^2]$ function is the \textbf{standard textbook model} (Balanis, Collin \& Zucker). Correctly produces wider H-plane beam and lower H-plane sidelobes than E-plane. \\
Singularity handling & A & L'H\^opital's rule correctly applied at $v = \pi/2$ where the denominator vanishes. \\
\texttt{horn\_aperture\_wavelengths} & B$-$ & Defined in \texttt{default\_params()} and prompted during \texttt{configure()}, but \textbf{not actually used} in \texttt{compute\_point\_gain()}. The value is collected but never wired to beamwidth derivation. \\
\texttt{ftb\_ratio\_db} & B & User knob. Real horn FTB depends on flare angle and waveguide modes. \\
Back lobe & B & Generic heuristic. Uses $0.3 \times n_\text{el}$ broadening factor (tightest of all types). \\
\texttt{n\_az}, \texttt{n\_el} & C+ & Only \texttt{n\_el} used for back lobe; \texttt{n\_az} is dead. \\
\bottomrule
\end{longtable}

\subsection{Shortcomings}

\begin{enumerate}[nosep]
    \item \textbf{\texttt{horn\_aperture\_wavelengths} is collected but unused.} The configure step prompts the user for aperture size and states ``beamwidth is derived from aperture size,'' but \texttt{compute\_point\_gain()} never references this value. This is a broken contract with the user.
    \item \textbf{No flare angle or waveguide geometry.} Horn beamwidth should be derivable from aperture dimensions $a \times b$ and flare semi-angles $\psi_E$, $\psi_H$.
    \item \textbf{No phase error (quadratic taper).} Real horns have a quadratic phase error across the aperture due to the spherical wavefront from the throat. This broadens the beam and raises sidelobes. The model assumes zero phase error (ideal horn).
    \item \textbf{No gain from geometry.} Standard gain horn gain is $G = \frac{4\pi}{\lambda^2} \cdot A_\text{eff}$, which is not computed.
    \item \textbf{No frequency scaling.} Horn beamwidth narrows with increasing frequency ($\text{BW} \propto \lambda / a$); this is not automatic.
\end{enumerate}

\subsection{Potential Improvements}

\begin{enumerate}[nosep]
    \item \textbf{Wire up \texttt{horn\_aperture\_wavelengths}}---the most impactful fix. Derive $\text{BW}_\text{az} \approx 67\lambda/a$ and $\text{BW}_\text{el} \approx 51\lambda/b$ from aperture dimensions.
    \item Add quadratic phase error parameter $s = a^2 / (8\lambda R)$ that modifies the pattern via the Fresnel integral formulation.
    \item Compute standard gain horn directivity from aperture area.
    \item Accept E-plane and H-plane aperture dimensions separately (rectangular horn is not always square).
\end{enumerate}


% =================================================================
\newpage
\section{Phased Array}
% =================================================================

\grade{afblue}{A$-$}

\subsection{Model Summary}

Implements the \textbf{Pattern Multiplication Theorem}:
$G_\text{total} = G_\text{element} \times \text{AF}$, where the Array Factor is:
\[
\text{AF}(\theta) = \frac{\left|\sum_{n=0}^{N-1} w_n \, e^{\,j n (kd\sin\theta + \beta)}\right|^2}
    {\left(\sum_{n} w_n\right)^2}
\]
with progressive phase steering $\beta = -kd\sin\theta_\text{steer}$.

Supports linear, planar (separable $\text{AF}_x \times \text{AF}_y$), and circular
array geometries.

\subsection{Parameter Assessment}

\begin{longtable}{@{}llp{6cm}@{}}
\toprule
\textbf{Parameter} & \textbf{Grade} & \textbf{Notes} \\
\midrule
\endhead
Array Factor & A+ & Exact wave superposition. Grating lobes, beam steering, and null positions are all physically correct. \\
Progressive phase steering & A+ & Standard formula $\beta = -kd\sin\theta_s$. Works for both linear and planar arrays. \\
Taylor weighting & B+ & Cosine-on-pedestal approximation, not exact Taylor. Produces reasonable sidelobe control but the exact sidelobe levels may differ from true Taylor by 1--2\,dB. \\
Mutual coupling & B & Exponential decay model $Z_{mn} = c_0 e^{-\alpha |m-n| d}$ with distance-dependent phase. Captures the trend but magnitude and phase are not derived from element geometry. Gaussian elimination solver is correct. \\
Element patterns & B$-$ & Three hardcoded models (Gaussian, cosine, patch) that are \textbf{not derived from element physics}. A ``patch'' element uses $\cos^{1.5}$ rolloff---empirical, not from cavity model. \\
RMS errors & A & Per-element Gaussian amplitude/phase perturbations with deterministic seed. Models real manufacturing tolerances correctly. \\
Circular array & A$-$ & Correct 3D phase computation with radial and parallel orientations. Arc-spacing approximation for coupling is reasonable. \\
Grating lobe warning & A & Correctly flags $d > \lambda/2$. \\
\texttt{ftb\_ratio\_db} & B & User knob. Real array FTB depends on element pattern and back-radiation, not an independent parameter. \\
Back lobe & B$-$ & Simplest back-lobe model of all types: just $\text{FTB} \times (1-\cos\phi)/2$ with no elevation shaping. \\
\bottomrule
\end{longtable}

\subsection{Shortcomings}

\begin{enumerate}[nosep]
    \item \textbf{Element patterns are empirical.} The three built-in element models are parametric fits, not derived from any physical element. A real array's performance depends critically on the element pattern, especially at wide scan angles.
    \item \textbf{Taylor weighting is approximate.} The cosine-on-pedestal formula does not exactly reproduce Taylor $\bar{n}$ weights. For precise sidelobe specifications, exact Taylor or Dolph--Chebyshev weights are needed.
    \item \textbf{Mutual coupling is frequency-independent.} The coupling matrix depends on element spacing in wavelengths but does not vary with frequency. Real coupling changes significantly across the band.
    \item \textbf{No scan impedance.} Active element impedance varies with scan angle (scan blindness at certain angles). This is not modelled.
    \item \textbf{No frequency scaling of spacing.} Element spacing is specified in wavelengths (fixed electrical spacing), not in physical metres. A real array with fixed physical spacing would see $d/\lambda$ increase with frequency.
    \item \textbf{Back lobe has no elevation taper.} Unlike all other antenna types, the array back lobe is flat in elevation.
    \item \textbf{Planar AF is separable only.} The product $\text{AF}_x \times \text{AF}_y$ is valid for rectangular lattices but not for triangular or hexagonal lattices.
\end{enumerate}

\subsection{Potential Improvements}

\begin{enumerate}[nosep]
    \item Allow any registered antenna type as the element pattern (already hinted at in the docstring but not implemented).
    \item Implement exact Taylor $\bar{n}$ weights and Dolph--Chebyshev weights.
    \item Add physical element spacing (metres) option with frequency-dependent $d/\lambda$ computation.
    \item Add scan impedance variation using an active-element-pattern model.
    \item Support triangular and hexagonal lattice geometries.
    \item Add elevation shaping to the back-lobe model for consistency with other antenna types.
\end{enumerate}


% =================================================================
\newpage
\section{LPDA (Log-Periodic Dipole Array)}
% =================================================================

\grade{aforange}{B+}

\subsection{Model Summary}

The LPDA uses an \textbf{elliptical Gaussian beam} model:
\[
G(\phi, \theta) = G_\text{peak} \cdot 0.5^{\,r^2} \cdot f_\text{sigmoid}(\cos\phi)
\]
where $r^2 = (\phi / \phi_\text{hw})^2 + (\theta / \theta_\text{hw})^2$ and
$f_\text{sigmoid}$ is a smooth front-hemisphere fade.

Back lobe: $G_\text{back} = G_\text{peak} \cdot \text{FTB}_\text{lin} \cdot
\frac{1 - \cos\phi}{2} \cdot |\cos\theta|^{0.5 n_\text{el}}$.

Total pattern: $G = \max(G_\text{front}, G_\text{back})$.

\subsection{Parameter Assessment}

\begin{longtable}{@{}llp{6cm}@{}}
\toprule
\textbf{Parameter} & \textbf{Grade} & \textbf{Notes} \\
\midrule
\endhead
Main beam (Gaussian) & A & Smooth elliptical iso-gain contours. Avoids rectangular artefacts of separable $\cos^n$ models. The $0.5^{r^2}$ formulation guarantees $-3$\,dB at the beamwidth contour. Good choice for a synthetic beam. \\
\texttt{az\_bw}, \texttt{el\_bw} & A & Direct, intuitive control. Correctly normalised via half-widths. \\
\texttt{ftb\_ratio\_db} & B+ & Clean $\text{dB} \to \text{linear}$ conversion. Static across frequency---real LPDA FTB degrades at band edges. \\
\texttt{ftb\_affects\_gain} & B$-$ & Rough energy conservation: $G' = G / (1 + \text{FTB}_\text{lin})$. Ignores sidelobe power. Better than nothing but undercounts total back-hemisphere radiation. \\
\texttt{sigmoid\_k} & B & Default $k = 40$ gives $\sim\!5^\circ$ transition at $\pm 90^\circ$---unrealistically sharp. An LPDA has a broader rolloff; $k = 10$--$15$ would be more representative. \\
Sidelobes (radial) & A$-$ & Elliptical rings via \texttt{compute\_sidelobe\_radial()} are a good model. Configurable first SLL, decay rate, and lobe count. Not LPDA-specific but reasonable. \\
Back lobe & A & Smooth $(1 - \cos\phi)/2$ weighting with no dead zones at $\pm 90^\circ$ endfire. Broader elevation ($0.5 \times n_\text{el}$) than front is physically reasonable. \\
\texttt{n\_az} & F & \textbf{Dead parameter.} Passed into \texttt{compute\_point\_gain()} but never referenced anywhere in the function body. \\
\texttt{n\_el} & C+ & Only used for back-lobe elevation broadening ($0.5 \times n_\text{el}$). Not used in the front-lobe Gaussian model. \\
\texttt{gain\_peak} & A & Correct role for a parameterised model. Passed through, optionally scaled by energy conservation. \\
\texttt{cfg} extensibility & A & \texttt{sigmoid\_k}, \texttt{ftb\_affects\_gain}, and all sidelobe parameters accessible via config dict without code changes. \\
\bottomrule
\end{longtable}

\subsection{Shortcomings}

\begin{enumerate}[nosep]
    \item \textbf{Generic beam shape.} The elliptical Gaussian is not LPDA-specific. The same model could represent a Yagi, log-spiral, or any other directional antenna. There is nothing in the pattern computation that reflects LPDA physics (active region, element superposition, feed network).
    \item \textbf{\texttt{n\_az} is a dead parameter.} It is computed by the engine from \texttt{az\_bw} via \texttt{bw\_to\_exponent()}, passed into the function, and then ignored. This is wasted computation and a confusing contract.
    \item \textbf{Sigmoid default is too sharp.} $k = 40$ creates a near-binary front/back division. An LPDA's radiation past $\pm 90^\circ$ tapers more gradually.
    \item \textbf{No frequency-dependent FTB.} LPDA front-to-back ratio typically degrades by 3--5\,dB at band edges where fewer elements are active. The model uses a single static value.
    \item \textbf{Energy conservation is crude.} The $1/(1 + \text{FTB}_\text{lin})$ correction does not account for sidelobe or wide-angle power.
    \item \textbf{Geometry module is disconnected.} \texttt{core/lpda\_geometry.py} computes Carrel-based element lengths, boom length, estimated directivity ($10 \log_{10}(N_\text{active} \cdot 2\sigma / (1-\tau))$), beamwidth ($70\lambda / L_\text{boom}$), and FTB ($8 + 40(\tau - 0.8)$). \textbf{None of these feed into the pattern computation.}
\end{enumerate}

\subsection{Potential Improvements}

\begin{enumerate}[nosep]
    \item \textbf{Wire geometry to pattern (highest impact).} Use \texttt{lpda\_geometry.py} outputs to auto-populate \texttt{az\_bw}, \texttt{el\_bw}, \texttt{gain\_peak}, and \texttt{ftb\_ratio\_db} when the user supplies $\tau$, $\sigma$, and frequency range. The Carrel estimates already exist---they just need to be connected.
    \item \textbf{Frequency-dependent FTB.} Model FTB degradation at band edges using a curve tied to the number of active elements at each frequency.
    \item \textbf{Reduce sigmoid default.} Change $k$ from $40$ to $12$--$15$ for a more realistic $\pm 30^\circ$ transition zone.
    \item \textbf{Remove or utilise \texttt{n\_az}.} Either delete it from the function signature or use it (e.g.\ as an alternative beam-shaping exponent when beamwidth is very wide).
    \item \textbf{Improve energy conservation.} Integrate sidelobe power into the conservation approximation, or replace with a proper pattern-integral normalisation.
    \item \textbf{Active-region beamwidth scaling.} Beamwidth should narrow at higher frequencies (more active elements, longer effective boom). This could be done via the existing engine frequency-dependent BW feature, but should be auto-configured from $\tau$ and $\sigma$ rather than requiring manual setup.
\end{enumerate}


% =================================================================
\newpage
\section{Shared Effects (\texttt{core/pattern\_math.py})}
% =================================================================

\grade{afblue}{B+}

These effects are applied by the engine to \textbf{all antenna types} after
\texttt{compute\_point\_gain()} returns. They are assessed independently of any
specific antenna.

\subsection{Effect Assessment}

\begin{longtable}{@{}lllp{5.5cm}@{}}
\toprule
\textbf{Effect} & \textbf{Grade} & \textbf{Basis} & \textbf{Notes} \\
\midrule
\endhead
Front-hemisphere sigmoid & B & Heuristic & Smooth $1/(1+e^{-k\cos\phi})$ transition. Eliminates heatmap artefacts. Not physics-derived but effective. \\
$\text{BW} \leftrightarrow n$ conversion & A & Exact math & $n = \ln 0.5 / \ln(\cos(\text{BW}/2))$ is the exact inverse. \\
Frequency-dependent BW & A$-$ & Aperture theory & Three modes (exponent, percentage, three-point). Exponent mode implements $\text{BW} \propto (f_\text{ref}/f)^\alpha$. Well-designed but requires manual feature configuration. \\
Sidelobe envelope & B+ & Empirical & Taylor-decay model with cosine-modulated lobes. Elliptical radial version produces realistic ring patterns. Not antenna-specific. \\
Pattern breakup & B & Synthetic & Deterministic sinusoidal ripple with onset angle. Visually realistic but not derived from any physical mechanism. \\
Ground reflection & A & Image theory & 2-ray model: $1 + \rho^2 + 2\rho\cos(\Delta\varphi + \pi)$. Correct interference fringe spacing. \\
Cross-polarisation & B & Empirical & Gaussian-envelope cross-pol with $\sin(2\arctan)$ diagonal factor. Produces the correct 45\textdegree\ peak pattern but levels are user-specified, not derived. \\
VSWR rolloff & B & Parametric & Cosine/linear taper at band edges. Models the trend but not derived from impedance mismatch. \\
Gain--BW coupling & A & Directivity law & $G \propto 1/(\text{BW}_\text{az} \cdot \text{BW}_\text{el})$. Correctly preserves the directivity relationship. \\
Asymmetry & B+ & Empirical & Pointing offset + deterministic pseudo-random perturbation. Models real installation misalignment. \\
\bottomrule
\end{longtable}

\subsection{Shortcomings}

\begin{enumerate}[nosep]
    \item \textbf{VSWR rolloff is not derived from impedance.} It applies a parametric gain penalty at band edges rather than computing $\text{VSWR} = (1 + |\Gamma|) / (1 - |\Gamma|)$ from actual antenna impedance.
    \item \textbf{Ground reflection uses a fixed $\rho$.} Real reflection coefficient depends on polarisation, frequency, grazing angle, and soil permittivity.
    \item \textbf{Cross-pol model has no polarisation basis.} There is no Jones matrix, Stokes parameter, or axial ratio computation. The model is a gain-envelope approximation only.
    \item \textbf{Sidelobe model is antenna-agnostic.} The same envelope is used for all antenna types despite very different sidelobe structures (e.g.\ dish Airy sidelobes vs.\ array grating lobes).
\end{enumerate}

\subsection{Potential Improvements}

\begin{enumerate}[nosep]
    \item Add angle-dependent and polarisation-dependent ground reflection coefficient (Fresnel equations).
    \item Compute VSWR from antenna impedance when impedance data is available.
    \item Allow per-antenna-type sidelobe models (the dish already has Jinc$^2$ sidelobes built into its front lobe; the engine sidelobe overlay may double-count).
    \item Add Ludwig-3 cross-pol definition for linearly polarised antennas.
\end{enumerate}


% =================================================================
\newpage
\section{Summary}
% =================================================================

\begin{center}
\begin{tabular}{@{}lcccccc@{}}
\toprule
\textbf{Antenna} & \textbf{Front Lobe} & \textbf{Freq.\ Scaling} & \textbf{Geom.\ $\to$ Pattern} & \textbf{Back Lobe} & \textbf{Params Used} & \textbf{Overall} \\
\midrule
Monopole & Exact formula   & Auto   & Full    & Omni (correct) & All   & \textbf{A+} \\
Dish     & Exact Airy      & Manual & Partial & Heuristic      & All   & \textbf{A}  \\
Panel    & Exact $\sinc^2$ & Manual & None    & Heuristic      & Mostly & \textbf{A}  \\
Horn     & Exact E/H       & Manual & \textbf{Broken}\footnotemark & Heuristic      & Mostly & \textbf{A}  \\
Array    & Exact AF        & Manual & Partial & Flat           & All   & \textbf{A$-$}  \\
LPDA     & Empirical       & Manual & \textbf{Disconnected} & Heuristic      & Mostly & \textbf{B+} \\
\bottomrule
\end{tabular}
\end{center}
\footnotetext{\texttt{horn\_aperture\_wavelengths} is prompted but never used in pattern computation.}

\bigskip

\subsection{Top Priority Improvements (Cross-Cutting)}

\begin{enumerate}
    \item \textbf{Wire LPDA geometry to pattern.} The Carrel equations in \texttt{lpda\_geometry.py} already compute beamwidth, directivity, and FTB estimates. Connecting them to the LPDA pattern generator would elevate it from B+ to A$-$ with no new physics required.

    \item \textbf{Wire horn aperture to beamwidth.} \texttt{horn\_aperture\_wavelengths} is collected from the user but ignored. This is a broken promise that should be fixed.

    \item \textbf{Add geometry-derived beamwidth to dish and panel.} Accepting aperture dimensions and deriving $\text{BW} = k\lambda/D$ would make frequency scaling automatic rather than requiring manual feature configuration.

    \item \textbf{Reduce sigmoid default from $k=40$ to $k=12$.} Applies to all directional antenna types. The current default produces an unrealistically sharp front/back boundary.

    \item \textbf{Clean up dead parameters.} \texttt{n\_az} is unused in LPDA, Panel, and Horn. Either utilise it or remove it from the call signature.
\end{enumerate}

\end{document}
